\subsection*{Theories}
\label{sec:axiom-systems-theories}

\begin{toolkitbox}{Theories Toolkit --- Quick Reference}
\small
\begin{tabular}{l l l}
\toprule
\textbf{Concept} & \textbf{Role} & \textbf{Detail} \\
\midrule
Theory & all consequences generated by axioms & \hyperref[def:axiom-systems-theory]{Definition} \\
Theorem & statement derivable from axioms & \hyperref[def:axiom-systems-theorem]{Definition} \\
Consequence & statement following from a theory & \hyperref[def:axiom-systems-consequence]{Definition} \\
Derivation & chain of justified steps & forward reference \\
Deductive closure & all derivable statements & \hyperref[def:deductive-closure]{Definition} \\
\bottomrule
\end{tabular}
\end{toolkitbox}

\begin{tcolorbox}[colback=structuretheorybox, colframe=structuretheoryborder, arc=2pt,
  left=8pt, right=8pt, top=6pt, bottom=6pt,
  title={\small\textbf{Theory Presentation: Arithmetic Consequences}},
  fonttitle=\small\bfseries]
Start with the arithmetic-style axiom system from the previous section. It
uses the arithmetic signature and includes assertions such as
\[
  \forall x\;(x+0=x),
\]
\[
  \forall x\;(S(x)\neq 0),
\]
\[
  \forall x\forall y\;
  (S(x)=S(y)\rightarrow x=y).
\]

Those axioms are the starting assertions. From them, together with other
arithmetic axioms and ordinary logical reasoning, additional statements may be
established. For instance, one may later prove new equations, cancellation
facts, uniqueness facts, and order facts that were not listed as axioms.

A theory contains the original axioms. A theory also contains the statements
derivable from those axioms. Thus the axioms are the generators, while the
theory is the full collection of what those generators imply.
\end{tcolorbox}

\begin{remark*}[Structural remarks]
\textbf{Axioms vs theory.}
Axioms are starting assertions. A theory includes all logical consequences of
those assertions. The distinction is the next step after the vocabulary and
assertion distinction: a signature supplies vocabulary, axioms supply adopted
assertions, and a theory contains everything derivable from those assertions.

\textbf{Finite axioms, infinite theory.}
A small axiom system can generate an enormous collection of consequences. The
listed axioms may be finite, but the statements derivable from them need not
be. This is why a theory is usually larger than the axiom system that presents
it.

\textbf{Forward roadmap.}
The next conceptual step is semantic:
\[
\begin{array}{c}
\text{Theory}\\
\Downarrow\\
\text{Models}
\end{array}
\]
A model will later be introduced as a mathematical structure in which the
statements of a theory hold.
\end{remark*}

\begin{definition}[Theory]
\label{def:axiom-systems-theory}
A \emph{theory} is a collection of statements in a language that contains the
axioms and all statements derivable from them.
\end{definition}

\begin{remark*}[Interpretation]
A theory is not merely the list of axioms. It is the closure of those axioms
under derivability: every statement that follows from the axioms belongs to
the theory. The axioms start the system; the theory records everything the
system proves.
\end{remark*}

\begin{dependencies}
\begin{itemize}
  \item \hyperref[def:language]{Language}
  \item \hyperref[def:axiom-system]{Axiom System}
\end{itemize}
\end{dependencies}

\begin{definition}[Theorem]
\label{def:axiom-systems-theorem}
A \emph{theorem} is a statement derivable from the axioms of a theory.
\end{definition}

\begin{remark*}[Interpretation]
Within an axiom system, a theorem is not an additional starting assumption. It
is a statement reached from the adopted axioms by reasoning. Axioms are given;
theorems are obtained.
\end{remark*}

\begin{dependencies}
\begin{itemize}
  \item \hyperref[def:axiom-systems-theory]{Theory}
\end{itemize}
\end{dependencies}

\begin{definition}[Consequence]
\label{def:axiom-systems-consequence}
A statement is a \emph{consequence} of a theory if it follows from the
statements of that theory.
\end{definition}

\begin{remark*}[Interpretation]
Consequence is the relation that connects a theory to what it implies. If the
axioms generate a theory, then consequences are the statements carried along
by those axioms and by the statements already obtained from them.
\end{remark*}

\begin{dependencies}
\begin{itemize}
  \item \hyperref[def:axiom-systems-theory]{Theory}
\end{itemize}
\end{dependencies}

\begin{definition}[Deductive Closure]
\label{def:deductive-closure}
The \emph{deductive closure} of an axiom system is the collection of all
statements derivable from that axiom system.
\end{definition}

\begin{remark*}[Interpretation]
Deductive closure is the operation that turns an axiom system into the full
theory it generates. Starting with the axioms, close under derivability; the
result is the theory determined by those starting assertions.
\end{remark*}

\begin{dependencies}
\begin{itemize}
  \item \hyperref[def:axiom-system]{Axiom System}
\end{itemize}
\end{dependencies}
